\documentclass[11pt,a4paper]{article}
\usepackage[T1]{fontenc}
\usepackage{lmodern,amsmath,amssymb,amsthm,microtype}
\usepackage[margin=27mm]{geometry}
\usepackage[colorlinks=true,linkcolor=blue,citecolor=blue,urlcolor=blue]{hyperref}
\setlength{\emergencystretch}{2em}
\newtheorem{theorem}{Theorem}[section]
\newtheorem{proposition}[theorem]{Proposition}
\newtheorem{lemma}[theorem]{Lemma}
\newcommand{\Q}{\mathbb Q}
\DeclareMathOperator{\Tr}{Tr}
\DeclareMathOperator{\Jac}{Jac}
\title{Arithmetic reductions for an actual-trace square equation}
\author{}\date{17 November 2026}
\begin{document}\maketitle
\begin{abstract}
For a fixed rational elliptic point tuple, the derivative detector's
actual-trace square equation factors through a curve in its physical
weight, of genus at most $36$. Rational points on that curve lift to
the source precisely through a degree-nine equation, a rational-square
condition, and the retained guards. An exact point tuple gives an
intermediate curve of genus $24$, a squared-slope quotient of genus
$224$, and a slope curve of genus $449$. The latter Jacobian is
isogenous over $\Q$ to the product of Jacobians of genera $224$
and $225$. Integer polynomial specifications and finite-field Bezout
certificates verify these genera. A separate exact slope has a soluble
split contact and a unit-weight local obstruction. No uniform rational
realization or defining-variable reduction is asserted.
\end{abstract}

\section{The actual weight curve}

Put $K=\Q(\sqrt2)$ and use the source curve and point
\[
 E:y^2=x^3-(3-\sqrt2)^2x,\qquad
 P=\left(\frac{57-28\sqrt2}{8},\frac{-660+417\sqrt2}{32}\right).
\]
Fix a guarded rational point tuple of the derivative source
\cite{Derivative,Criterion}. Its six tests $C_i\in K$ are actual
constants; the trace is not an independent parameter. For a variable
$u$ write
\begin{equation}\label{eq:actual}
 D_i=\frac{C_i^2}{1+uC_i^2},\quad
 W=\prod_{i=1}^6(1+D_i^2/u),\quad
 V=\Tr(W^{-1}),\quad
 \mathcal C(u)=V^2+\frac{u}{(1+u)^2}(2-V).
\end{equation}
In the source the first four tests are nonzero. The two comparison
tests may vanish, and a terminal ordinate may be zero.

\begin{proposition}\label{prop:small}
For any such fixed tuple, the normalization of
$Y^2=\mathcal C(u)$ is geometrically integral and has genus at
most $36$. It has a rational branch point at $u=0$ and two rational
points over $u=\infty$, outside the guarded affine source.
\end{proposition}
\begin{proof}
For each test put
\[
 L_i(u)=u(1+uC_i^2)^2+C_i^4,\qquad
 S_i(u)=u(1+uC_i^2)^2.
\]
Both products $L=\prod_i L_i$ and $S=\prod_i S_i$ have
degree at most $18$ over $K$. Thus
\[
 V=\frac{S\overline L+\overline S L}{L\overline L}
      =\frac{N(u)}{D(u)},\qquad N,D\in\Q[u],\quad
                  \deg N,\deg D\le36.
\]
Common factors may be cancelled; the original guards are retained
independently of that cancellation. Clearing the square denominator
gives the polynomial model
\begin{equation}\label{eq:weight-polynomial}
 Y_0^2=\Pi(u):=N(u)^2(1+u)^2+uD(u)(2D(u)-N(u)),
 \qquad Y_0=(1+u)D(u)Y.
\end{equation}
Its degree is at most $74$. Removing square factors gives a
hyperelliptic normalization of genus at most $36$.

Every nonzero test contributes a simple zero to $W^{-1}$ at
$u=0$ in both embeddings. Since at least one test is nonzero,
$V=O(u)$, and $\mathcal C(u)=2u+O(u^2)$. The function is
therefore not a square even over $\overline\Q(u)$. This also
gives the rational branch point. At infinity,
$D_i=u^{-1}+O(u^{-2})$ for nonzero tests, while zero tests
contribute the factor one. Hence $V=2+O(u^{-3})$ and
$\mathcal C(u)=4+O(u^{-3})$, giving two rational points there.
\end{proof}

The genus statement concerns the smooth normalization. If a polynomial
model in \eqref{eq:weight-polynomial} is singular, its finitely many
rational singular points must also be tested directly in the affine
equation; a rational singular point need not lift rationally to its
normalization. This includes permitted zero terminal ordinates.

\section{An exact rational point tuple}

Take
\begin{equation}\label{eq:fixture}
 X=1,\qquad B=3,\qquad A=P,\qquad H=U=5P.
\end{equation}
Here $B=3$ is an algebraic fixture. No claim is made that it exceeds
the ineffective defining-shift threshold of the source theorem.
The identity
\[
                  2(1+\sqrt2,2\sqrt2)=P
\]
places the three point roles in $2E(K)$. The points $P,4P,5P$
have nonzero finite abscissas and ordinates, and $x(5P)\ne x(P)$;
these assertions are checked exactly by the accompanying program.

Let $f(Q)=1/x(Q)$ and put
\[
 c=f(5P),\qquad d=\overline{f(4P)}.
\]
The six source tests are exactly
\[
                    (C_1,\ldots,C_6)=(c,c,d,d,0,0).
\]
Indeed both formal quotients are one, since $U=H$ and $U-A=H-P$.
Writing $\alpha=c^2$ and $\beta=d^2$, define
\[
 \begin{aligned}
 L_\alpha&=u(1+u\alpha)^2+\alpha^2,
       &L_\beta&=u(1+u\beta)^2+\beta^2,\\
 Q&=L_\alpha L_\beta,
       &S&=u^2(1+u\alpha)^2(1+u\beta)^2.
 \end{aligned}
\]
Then $W=(Q/S)^2$. Consequently the actual trace is
\begin{equation}\label{eq:trace-fixture}
 V=\frac{(S\overline Q)^2+(\overline S Q)^2}
          {(Q\overline Q)^2}=\frac{N(u)}{D(u)}.
\end{equation}
The last pair is made integral by multiplying numerator and denominator
by the same common denominator and dividing their common integer
content. This specifies $N,D$ exactly; their ascending coefficient
lists are also included in the certificate.

\begin{theorem}\label{thm:fixture}
For \eqref{eq:fixture}, $N$ and $D$ are coprime of degree $24$,
with $\operatorname{lc}(N)=2\operatorname{lc}(D)$. The polynomial
$\Pi$ of \eqref{eq:weight-polynomial} is squarefree of degree $50$.
Thus the actual intermediate curve has genus $24$.
\end{theorem}
\begin{proof}
The exact polynomial construction verifies the degrees and leading
coefficients. Reduction modulo $101$ preserves the degrees and
certifies $\gcd(N,D)=1$. The certificate supplies polynomials
$e,f\in\mathbb F_{101}[u]$ satisfying
\[
                     e\Pi+f\Pi'=1.
\]
Thus $\Pi$ is squarefree over $\Q$, giving genus $(50-2)/2=24$.
As an independent arithmetic check of the trace formula, exact field
evaluation of \eqref{eq:actual} at $u=1,2,3$ agrees with $N/D$.
\end{proof}

\section{The physical map and its two quotients}

For \eqref{eq:fixture}, the scale is
\[
 a=\left(4^4+(57/4)^4\right)^4
          =\left(\frac{10621537}{256}\right)^4>10.
\]
Throughout this section $t=z^2$; it is not the source expression
$a^2-z^2$. For
$\mathcal S=\{-10,-8,-5,1,2,3,4,6,7\}$ put
\[
 F(a,t)=\prod_{c\in\mathcal S}((a+c)^2-t),\qquad
              \frac{\mathsf A(t)}{\mathsf B(t)}=-18\frac{F}{F_a},
\]
where $F_a$ is differentiated before $a$ is specialized.
The integer pair $\mathsf A,\mathsf B$ is obtained by the same
common-denominator and common-content convention as above. It is
coprime, with degrees $9$ and $8$. Its coefficient lists are in the
certificate. Define the polynomial by homogeneous composition:
\begin{equation}\label{eq:composition}
 \mathcal P(t)=\mathsf B(t)^{50}
                 \Pi\!\left(\frac{\mathsf A(t)}{\mathsf B(t)}\right).
\end{equation}
This is an exact integer polynomial specification without a printed
degree-$450$ expansion. On $\mathsf B(t)\ne0$ put
$Y_1=\mathsf B(t)^{25}Y_0$; this identifies the displayed
composition curves with the corresponding actual equations.

\begin{theorem}\label{thm:genera}
The polynomial $\mathcal P$ is squarefree of degree $450$, and
$\mathcal P(0)\ne0$. The normalized curves
\[
 \mathcal C_t:Y_1^2=\mathcal P(t),\quad
 \mathcal C_t':Y_2^2=t\mathcal P(t),\quad
 \mathcal C_z:Y_1^2=\mathcal P(z^2)
\]
have genera $224,225,449$, respectively. There is an isogeny over
$\Q$
\[
              \Jac(\mathcal C_z)\longrightarrow
              \Jac(\mathcal C_t)\times\Jac(\mathcal C_t').
\]
\end{theorem}
\begin{proof}
The degree follows from $\deg\Pi=50$, $\deg\mathsf A=9$ and
$\deg\mathsf B=8$. Modulo $101$ this degree is preserved,
$\mathcal P(0)\ne0$, and a second recorded Bezout identity gives
$e_t\mathcal P+f_t\mathcal P'=1$. It follows that both
$t\mathcal P(t)$ and $\mathcal P(z^2)$ are squarefree. The
hyperelliptic genus formulas give the stated values.

The two maps are
\[
 (z,Y_1)\longmapsto(z^2,Y_1),\qquad
 (z,Y_1)\longmapsto(z^2,zY_1).
\]
Their pullbacks send the regular differential bases to
\[
 2z^{2i+1}\frac{dz}{Y_1}\quad(0\le i\le223),\qquad
 2z^{2j}\frac{dz}{Y_1}\quad(0\le j\le224).
\]
Together these form the full $449$-dimensional space of regular
differentials on $\mathcal C_z$. The product of the induced
Jacobian pushforwards therefore has invertible cotangent map.
In characteristic zero it is surjective with finite kernel, hence
an isogeny; see the basic Jacobian and isogeny constructions in
\cite{Milne}.
\end{proof}

\section{Exact rational lifting}

\begin{proposition}\label{prop:lift}
A finite rational point $(u,Y_0)$ on $Y_0^2=\Pi(u)$ gives an
actual rational square witness for the fixed tuple precisely when
there is $t\in\Q$ such that
\begin{equation}\label{eq:lift}
 u\mathsf B(t)-\mathsf A(t)=0,\qquad
 \mathsf B(t)\ne0,\qquad t=z^2\text{ for some }z\in\Q,
\end{equation}
and all retained source guards hold. In particular these include
\[
 u\ne0,-1,\qquad N(u)D(u)\ne0,\qquad a^2-t\ne0,
\]
the balanced denominator norms, both slots of $W$, and the original
point and difference guards. For a prepared witness one additionally
retains every prescribed fixed-place, selected-contact, real and
whole-character condition of \cite{Criterion}.
The equation in $t$ has degree exactly nine.
\end{proposition}
\begin{proof}
The finite weight $u$ is physical exactly when the first two
conditions of \eqref{eq:lift} hold. Its polynomial has degree nine
because $\deg\mathsf A=9>\deg\mathsf B$. The last condition
is precisely rational recoverability of the slope. Under the guards,
\[
                   Y=\frac{Y_0}{(1+u)D(u)}
\]
then satisfies the actual square equation, and every actual witness
arises this way. Zero $Y_0$ is allowed. Extra preparation conditions
are additional tests, not consequences of rationality on a quotient.
\end{proof}

Equivalently, rational points $(t,Y_1)$ and $(t,Y_2)$ on the two
squared-slope quotients with the same $t$ and $Y_1\ne0$ give
$z=Y_2/Y_1\in\Q$. When $Y_1=0$, both quotient ordinates
vanish and the separate rational-square test on $t$ is still needed.
The known rational points above $u=0$ and $u=\infty$ violate the
affine source guards. Thus the reductions do not by themselves
produce an admissible rational point.

\section{An exact local comparison at one slope}

Keep \eqref{eq:fixture} and set $z=a+12$. Then $z>a+10$ and
$u>0$: every summand in
\[
                     \frac{F_a}{F}
               =\sum_{c\in\mathcal S}\frac{2(a+c)}{(a+c)^2-z^2}
\]
is negative. All balanced denominators and factors are positive in
both real embeddings, so $0<V<2$ and $\mathcal C>0$.

At $17$, the contact $z\equiv a-5$ is isolated and
$v_{17}(u)=1$, with $u/17\equiv1\pmod{17}$. In the slot
$\sqrt2\equiv244\pmod{17^2}$, both $c$ and $d$ vanish modulo
$17^2$; in the other slot they are $2$ and $1$ modulo $17$.
Consequently $V\equiv1$ and $\mathcal C\equiv1\pmod{17}$,
so the actual radicand is a nonzero square in $\Q_{17}$.

At $31$, exact reduction of the same rational functions gives
\[
       u\equiv9,\qquad V\equiv22,\qquad
       \frac{u}{(1+u)^2}\equiv19,\qquad
                       \mathcal C\equiv11\pmod{31}.
\]
The balanced-factor and trace expressions have unit denominators
after reduction, while $11^{15}\equiv-1\pmod{31}$.
Thus this actual positive radicand is not a rational square. These
calculations verify the real condition and a single split contact;
they do not verify the retained preparation at every other place or
point-algebra character. Both $17$ and $31$ belong to the retained
fixed exceptional set. In particular this is not a counterexample
to rational realization after full preparation.

\section{Certificate scope}

The portable program uses exact rational arithmetic in $K$, exact
polynomial arithmetic, and the displayed finite-field Bezout
identities. It verifies the source point additions, trace identity,
coprimality and genera, and the two local calculations. The degree-$450$
polynomial is kept in homogeneous-composition form. The computations
do not enumerate rational points on the genus-$24$ curve and do not
show that a point surviving the lifting and preparation tests exists.

\begin{thebibliography}{9}
\bibitem{Derivative}\emph{An eleven-variable construction under rational $abc$},
\nolinkurl{papers/conditional_few_variables/derivative_detector11_20261117_v63.tex}.
\bibitem{Criterion}\emph{One actual-trace square in a derivative detector},
\nolinkurl{papers/pole_parameters/single_trace_derivative_criterion_20261117_v64.tex}.
\bibitem{Milne}J. S. Milne, \emph{Abelian Varieties}, version 2.0,
2008, Chapters I and III. \url{https://www.jmilne.org/math/CourseNotes/AV.pdf}.
\end{thebibliography}
\end{document}
